\documentclass[11pt, a4paper]{article}

\usepackage[T1]{fontenc}
\usepackage[utf8]{inputenc}
\usepackage{tgpagella}
\usepackage{microtype}
\usepackage[spanish, es-noshorthands]{babel}

\usepackage{amsmath, amssymb}
\usepackage[table, xcdraw]{xcolor}
\usepackage{booktabs}
\usepackage[most]{tcolorbox}
\usepackage[margin=2.5cm]{geometry}
\usepackage{fancyhdr}
\usepackage{pgfplots}
\pgfplotsset{compat=1.18}

\pagestyle{fancy}
\fancyhf{}
\setlength{\headheight}{16pt}
\lhead{\textbf{UCB Sede Tarija} | Ing. Mecatrónica}
\chead{Worksheet: Trayectorias LSPB}
\rhead{Semana 9 - Sesión 1}
\lfoot{Prof: M.Sc. Ing. Bernardo Quiroga Turdera}
\rfoot{Página \thepage}

\definecolor{ucbblue}{RGB}{0, 51, 102}

\begin{document}

\begin{center}
    \Large\textbf{\textcolor{ucbblue}{Hoja de Trabajo: Segmentos Lineales con Mezclas Parabólicas (LSPB)}}\\[0.2cm]
\end{center}

\vspace{0.2cm}

\section*{A. Análisis del Problema Guiado[cite: 13]}
Se desea planificar un movimiento articular con los siguientes parámetros[cite: 13]:
\begin{equation*}
    q_i = 0, \quad q_f = 2\text{ rad}, \quad t_f = 3\text{ s}, \quad \ddot{q}_c = 1\text{ rad/s}^2
\end{equation*}

\textbf{1. Factibilidad:} Calcule la aceleración mínima requerida $|\ddot{q}_c|_{\min} = \frac{4|q_f - q_i|}{t_f^2}$ y justifique si es factible construir un perfil trapezoidal con la $\ddot{q}_c$ dada.
\vspace{2cm}

\textbf{2. Tiempo de Mezcla:} Calcule el tiempo de aceleración $t_c$ usando:
\begin{equation*}
    t_c = \frac{t_f}{2} - \frac{1}{2}\sqrt{t_f^2 - \frac{4(q_f - q_i)}{\ddot{q}_c}} \text{[cite: 13]}
\end{equation*}
\vspace{2cm}

\textbf{3. Velocidad de Crucero:} Calcule $\dot{q}_c = \ddot{q}_c t_c$.
\vspace{1.5cm}

\section*{B. Gráficos Alineados y Continuidad[cite: 13]}
Los siguientes gráficos muestran los perfiles de posición, velocidad y aceleración para el problema anterior.
\vspace{0.4cm}

\begin{center}
\begin{tikzpicture}
    \begin{axis}[
        width=10cm, height=4cm,
        ylabel={$q(t)$ [rad]},
        xmin=0, xmax=3, ymin=0, ymax=2,
        xtick=\empty,
        grid=major
    ]
    \addplot[ucbblue, thick, domain=0:1] {0.5*x^2};
    \addplot[ucbblue, thick, domain=1:2] {x-0.5};
    \addplot[ucbblue, thick, domain=2:3] {2 - 0.5*(3-x)^2};
    \draw[dashed, red] (axis cs:1,0) -- (axis cs:1,2);
    \draw[dashed, red] (axis cs:2,0) -- (axis cs:2,2);
    \end{axis}
\end{tikzpicture}

\begin{tikzpicture}
    \begin{axis}[
        width=10cm, height=4cm,
        ylabel={$\dot{q}(t)$ [rad/s]},
        xmin=0, xmax=3, ymin=0, ymax=1.2,
        xtick=\empty,
        grid=major
    ]
    \addplot[ucbblue, thick, domain=0:1] {x};
    \addplot[ucbblue, thick, domain=1:2] {1};
    \addplot[ucbblue, thick, domain=2:3] {3-x};
    \draw[dashed, red] (axis cs:1,0) -- (axis cs:1,1.2);
    \draw[dashed, red] (axis cs:2,0) -- (axis cs:2,1.2);
    \end{axis}
\end{tikzpicture}

\begin{tikzpicture}
    \begin{axis}[
        width=10cm, height=4cm,
        xlabel={$t$ [s]}, ylabel={$\ddot{q}(t)$ [rad/s$^2$]},
        xmin=0, xmax=3, ymin=-1.2, ymax=1.2,
        xtick={0,1,2,3},
        grid=major
    ]
    \addplot[ucbblue, thick, domain=0:0.99] {1};
    \addplot[ucbblue, thick, domain=1.01:1.99] {0};
    \addplot[ucbblue, thick, domain=2.01:3] {-1};
    \draw[dashed, red] (axis cs:1,-1.2) -- (axis cs:1,1.2);
    \draw[dashed, red] (axis cs:2,-1.2) -- (axis cs:2,1.2);
    \end{axis}
\end{tikzpicture}
\end{center}

\textbf{1.} Identifique y marque en los gráficos las zonas correspondientes a las fases de \textit{aceleración}, \textit{velocidad constante} y \textit{desaceleración}[cite: 13]. \\
\textbf{2.} Analice la continuidad: ¿Existen saltos instantáneos (discontinuidades) en la gráfica de $\dot{q}(t)$? ¿Y en la gráfica de $\ddot{q}(t)$? ¿Por qué es físicamente aceptable (en modelos básicos) un salto en la aceleración pero no en la velocidad?[cite: 13]
\vspace{2cm}

\end{document}
